\chapter{Signal Timing}
\label{ch:signal-timing}

This chapter gives the timing of the pins of the device (Chapter~\ref{ch:package}): the bus cycles in which the CPU
reads and writes it, the interrupt output, reset, and the serial output. Chapter~\ref{ch:timing} gives the timing
inside the device, in the same units.

\section{Reading the Figures}

The figures follow the conventions of the bus timing diagrams of the \hwmanual{}. Each labeled interval is one
\texttt{CLKI} cycle, seen between rising edges. Waveform rows show physical levels: for an active-low signal,
\texttt{L} means asserted. Bands show a value; an empty band carries no defined value, and a three-state signal that is
not driven is drawn as a flat line at mid-level. Electrical delays are not specified, so the figures show logical
timing only.

\section{Bus Cycles}
\label{sec:signal-bus}

The device answers every cycle in which \activelow{CS} and \activelow{BAS} are both asserted. How soon it asserts
\activelow{BACK} depends on the address: a register is answered in the cycle itself, and a word of the audio memory
window when the device has made the access.

\subsection{Register Accesses}

A register read or write is completed in the cycle in which it is made: \activelow{BACK} is asserted from the start of
the cycle, so the CPU accepts the cycle at the next rising edge and does not wait. On a write the register holds the
new value from that edge, and a read in the next cycle returns it. A write to a reserved offset is acknowledged in the
same way and changes nothing.

\begin{figure}[H]
	\centering
	\small
	\begin{tikztimingtable}[sirc timing diagram]
		CLKI             & L 2{0.5H 0.5L}              \\
		Phase            & D{access} D{next} D{}       \\
		\activelow{CS}   & L H H                       \\
		\activelow{BAS}  & L H H                       \\
		BRW              & L H H                       \\
		A0--A15          & D{offset} D{} D{}           \\
		D0--D15          & D{data} Z Z                 \\
		\activelow{BACK} & L H H                       \\
		Register         & D{old} 2D{new}              \\
	\end{tikztimingtable}
	\caption{Register Write Cycle}
	\label{fig:signal-write}
\end{figure}

On a read the device drives D0--D15 with the register's value from the start of the cycle to its end, and releases the
bus after it. A read changes nothing in the device.

\subsection{Audio Memory Window Accesses}

An access to the window waits for a CPU slot of the memory schedule (Section~\ref{sec:cpu-slots}). The CPU holds its
signals and the device holds \activelow{BACK} deasserted until the access is made, then asserts \activelow{BACK} for
the cycle in which a read puts its data on D0--D15, or a write takes its data. A window access is acknowledged in the
second to fifth cycle of the access, whatever the voices are doing (Section~\ref{sec:cpu-slots}).
Figure~\ref{fig:signal-window-read} shows a read of the slowest case, five cycles.

\begin{figure}[H]
	\centering
	\small
	\begin{tikztimingtable}[sirc timing diagram]
		CLKI             & L 5{0.5H 0.5L}                    \\
		Phase            & D{access} 3D{wait} D{done} D{}    \\
		\activelow{CS}   & 5L H                              \\
		\activelow{BAS}  & 5L H                              \\
		BRW              & 6H                                \\
		A0--A15          & 5D{offset} D{}                    \\
		D0--D15          & 4Z D{data} Z                      \\
		\activelow{BACK} & 4H L H                            \\
	\end{tikztimingtable}
	\caption{Audio Memory Window Read Cycle with Waits}
	\label{fig:signal-window-read}
\end{figure}

\section{Interrupt Output}

\activelow{IRQ3} is asserted while any \texttt{SEQ\_DONE} flag is set and its interrupt is enabled, and released when no
such flag remains (Section~\ref{sec:interrupts}). It is level-sensitive; no edge is produced. Its delay from the change
of the flags is not yet specified.

\section{Reset}

While \activelow{RST} is asserted the device holds the reset state of Section~\ref{sec:reset}. The converter output is
digital zero from reset until the program sets \texttt{AUDIO\_EN}. No interrupt is raised by reset. The minimum width
of the reset pulse is not yet specified.

\section{Serial Output}
\label{sec:serial-output}

The device sends one stereo sample every 336 \texttt{CLKI} cycles. The frame has 32 bit clocks, \texttt{BCLK}, in two
slots of 16, one for each channel. The periods of \texttt{BCLK} are 10 and 11 \texttt{CLKI} cycles alternately, the
first of the frame being 10, so the frame is $16 \times 10 + 16 \times 11 = 336$ cycles. The pattern is the same in
every frame, so every edge of \texttt{BCLK} falls at the same cycle of the frame each time. \texttt{WS} is low during
the left slot and high during the right. The duty cycle of \texttt{BCLK} is not yet specified.

\begin{table}[H]
	\centering
	\begin{tabular}{llp{6.4cm}}
		\toprule
		\textbf{\texttt{BCLK} periods of the frame} & \textbf{\texttt{WS}} & \textbf{\texttt{SD}}                          \\
		\midrule
		0                                           & Low                  & Bit 0 of the right word of the previous frame \\
		1--15                                       & Low                  & Bits 15 to 1 of the left word                 \\
		16                                          & High                 & Bit 0 of the left word                        \\
		17--31                                      & High                 & Bits 15 to 1 of the right word                \\
		\bottomrule
	\end{tabular}
	\caption{Serial Frame Layout}
	\label{tab:serial-frame}
\end{table}

The format is that of a standard two-channel serial audio interface. \texttt{WS} changes one bit clock before the most
significant bit of the word it announces, so the slot is as wide as the word and the last bit of each word is sent in
the first period of the next slot, after \texttt{WS} has changed. The 16 data bits are the output sample as a signed
two's-complement number, most significant first. A slot carries nothing else. The alignment of the edges of \texttt{WS}
and \texttt{SD} to \texttt{BCLK}, their set-up and hold times, and the \texttt{SCK} frequency and phase, are not yet
specified.
